\documentclass[a4paper, 10pt]{article}
\usepackage[unicode]{hyperref}
\usepackage[utf8x]{inputenc}
\usepackage[IL2]{fontenc}
\usepackage[slovak]{babel}
\usepackage[left=2cm, top=3cm, text={17cm, 24cm}]{geometry}
\usepackage{verbatim}
\usepackage{enumitem}
\usepackage{graphicx}
\usepackage{listings} % vkladanie obrazkov
\usepackage{color}
\usepackage{xcolor}
\usepackage[czech, boxed]{algorithm2e}


\begin{document}
%%%%%%%%%%%%%%%%%%%%%%%%%%%%%%%% Titulná stránka %%%%%%%%%%%%%%%%%%%%%%%%%%%%%%%%
    \begin{titlepage}
        \begin{center}
            \includegraphics[width=0.77\linewidth]{src/FIT_logo} \\
            \vspace{\stretch{0.382}}
            \Huge{Modelování a simulace -- IMS} \\
            \huge{Okruh 12. -- SHO vo výrobe} \\
            \Large{Dokumentácia}
            \vspace{\stretch{0.618}}
        \end{center}

        \begin{flushleft}
            \Large{Lapeš Zdeněk (xlapes02)} \\
            \Large{Bínovský Andrej (xbinov00)}
        \end{flushleft}
        \vspace{-12mm}
        \hfill\Large{\today}
    \end{titlepage}


%%%%%%%%%%%%%%%%%%%%%%%%%%%%%%%% Obsah %%%%%%%%%%%%%%%%%%%%%%%%%%%%%%%%
    \tableofcontents
    \newpage

%%%%%%%%%%%%%%%%%%%%%%%%%%%%%%%% Úvod %%%%%%%%%%%%%%%%%%%%%%%%%%%%%%%%


    \section{Úvod}
    V tejto práci je riešená implementácia procesu výroby chleba, ktorá je použitá na zostavenie modelu testujúceho
    najlepšie možné nakonfigurovanie množstva strojov, pracovníkov a miestností výroby pre rôzne scenáre požiadaviek
    množstva chlebov v daný deň. Danú prácu vypracovali študenti Zdenek Lapeš a Andrej Bínovský z Fakulty informačných
    technológií VUT v Brne.

    \subsection{Čerpanie a konzultácia dát}
    V rámci tejto práce boli dáta použité a konzultované z praxe, ktoré boli získané z výrobného procesu chleba
    z Brnenskej pekárne \texttt{Crocus}\cite{Crocus}. Pre dosianutie čo najvalidnejšieho modelu procesu výroby chleba sme
    spomenutú pekáreň navštívili a získali informácie o problematike každého procesu výroby.

    \subsection{Overovanie validity dát}
    Po celý čas modelovania boli dáta overované a validované na základe komunikácie majitela pekárne. Hned zo
    začiatku sme si namodelovali a porovnali presnú konfiguráciu spomenutej pekárne. Čo nám docielilo overenie
    presnosti modelu na základe porovnania skutočného času výroby s časom výroby chleba v modeli.


%%%%%%%%%%%%%%%%%%%%%%%%%%%%%%%% Rozbor témy a použitých metód/technológií %%%%%%%%%%%%%%%%%%%%%%%%%%%%%%%%
    \section {Rozbor témy}\label{sec:rozbor-temy}
    Postup výroby chleba sa skladá z viacerych procesov, ktoré sú navzájom závislé:
    \begin{itemize}
        \item \textbf{Výroba cesta} -- Proces výroby cesta spočíva zo zmiešania všetkých surovín do jedného. Počet
        a dostupnosť surovín sa v modeli neberie k úvahe. Jeden proces výroby cesta reprezentuje jeden mixér a
        výsledokm je \textbf{140kg cesta}. Proces trvá \textbf{10 minút} \textpm \textbf{ 2 minúty}.

        \item \textbf{Krájanie cesta na bochníky} -- Krájanie cesta na bochníky je proces, ktorý sa vykonáva pracovníkom
        ručne. V modeli teda platí, že jeden pracovník sa rovná jednému stolu na krájanie. Vstupom procesu je
        \textbf{1kg cesta} a výstupom je \textbf{1 bochník}. Proces trvá \textbf{30 sekúnd} \textpm \textbf{ 5 sekúnd}.

        \item \textbf{Fermentácia bochníkov} -- Po nakrájaní sa bochníky ukladajú na plech do vozíka.
        Jeden vozík obsahuje miesto na \textbf{70 bochníkov}. Po naplnení sa vozík odvezie do fermentačnej miestnosti
        na čas  \textbf{20 minút} \textpm \textbf{ 2 minúty}.

        \item \textbf{Pečenie chleba} -- Po fermentácii sa bochníky na vozíkoch uložia do pece, kde sa pečú.
        Jedna pec má kapacitu \textbf{1 vozíku}, teda \textbf{70 bochníkov}. Proces trvá \textbf{30 minút} \textpm \textbf{ 2 minúty}.

        \item \textbf{Balenie chleba} -- Po upečení pracovníci roztrieda chleba do bedien. Balenie jedného chleba
        praconíkom trvá \textbf{10 sekúnd} \textpm \textbf{ 3 sekundy}.
    \end{itemize}

    \subsection{Použité metódy, postupy a technológie}
    Na implementáciu modelu bol použitý programovací jazyk \texttt{C++} za podpory simulačnej knižnice \texttt{SimLib}.
    Hlavným dôvodom použitia jazyka \texttt{C++} je využitie naimplementovanej
    knižnice \texttt{SimLib}. Knižnicu sme si vybrali z dôvodu modelovania procesov výroby na základe petriho siete, ktorá
    toto modelovanie podporuje. Príklady a správne použitie knižnice sme čerpali z pradnášok predmetu \texttt{IMS}\cite{IMS_slides}.
    Na kompiláciu sme využili nástroj \texttt{CMake}. Na vývoj bol zvolený operačný sýstém
    \texttt{Linux Ubuntu}, ktorý bežal vo virtualizovanom prostredí nástroja \texttt{Docker}.

    \subsection{Pôvod metód a technológií}
    \begin{itemize}
        \item \textbf{C++} -- Verzia \texttt{C++20}: \texttt{https://en.cppreference.com/w/cpp/20}.
        \item \textbf{SimLib} -- Verzia \texttt{3.09}\cite{simblib}. Autori knižnice sú Petr Peringer,
        David Leska a David Martinek.
        \item \textbf{CMAKE} -- Verzia \texttt{3.5}: \texttt{https://cmake.org/}.
        \item \textbf{Docker} -- Verzia \texttt{20.10.17}: \texttt{https://www.docker.com/}.
        \item \textbf{Linux Ubuntu} -- Verzia \texttt{20.04}: \texttt{https://ubuntu.com/}.
    \end{itemize}


%%%%%%%%%%%%%%%%%%%%%%%%%%%%%%%% Implementácia programu %%%%%%%%%%%%%%%%%%%%%%%%%%%%%%%%
    \section {Koncepcia modelu}
    Konceptuálny model pekárne je zobrazený pomocou petriho siete %\ref{fig:conceptual_model}.
    V modeli je zobrazený zjednodušený proces výroby chleba. Napriek zjednodušeniu modelu je zachovaná korektnosť simulácie.
    \begin{itemize}
        \item \textbf{Výroba cesta} -- Pri prvom procese je zanedbateľná dostupnosť a naskladnovaie surovín,
        pretože to zaobstaráva externá firma. Výstup a čas procesu je totožný s realitou.

        \item \textbf{Krájanie cesta na bochníky} -- Pri druhom procese sme reálne dáta dekomponovali. Získané
        dáta z pekárne reprezentovali čas nakrájania 70 bochníkov (zo 70 kg cesta) na jeden vozík pracovníkmi za čas \textbf{10 min} \textpm \textbf{ 2 minúty}.
        Pre systém sme zvolili proces krájania 1kg cesta na 1 bochník, 1 pracovníkom za čas \textbf{30 sekúnd} \textpm \textbf{ 5 sekúnd}.

        \item \textbf{Fermentácia bochníkov} -- Pri procese fermentácie sme zachovali vstupy, výstupy a časy na základe reálnych dát.

        \item \textbf{Pečenie chleba} -- Pri procese pečenia sme zachovali vstupy, výstupy a časy na základe reálnych dát.

        \item \textbf{Balenie chleba} -- Proces balanie chleba reprezentuje viacero podprocesov ako napríklad príprava a čistenie bedien.
        Balenie spočíva v roztriedení chleba do bedien a ich následne uloženie do skladu. Reálny proces bol meraný na
        zabalenie jedného vozíku 70 chlebov do bedien za čas \textbf{10 min} \textpm \textbf{ 2 minúty}. Jedna bedna môže obsahovať 3 až 4 chleby. Proces v modeli sme dekomponovali
        na zabalenie jedného chleba jedným pracovníkom za čas \textbf{10 sekúnd} \textpm \textbf{ 3 sekundy}.
    \end{itemize}

    \section{Architektúra simulačného modelu}
    Simulačný model\cite[slide 44]{IMS_slides} je možné spustiť v 3 režimoch. Nastavenie týchto
    režimov je určené konštantami v súbore \texttt{src/macros.h}.
    Každý proces ma nasledujúcu implementáciu funkcie \texttt{Behavior()}:
    Na začiatku sa spustia všetky \texttt{mixer process}, poprípade se vytvori fronta na \texttt{mixer process}.

    \begin{algorithm}[ht]
        \SetKw{Break}{break}

        Zahájenie order process\;
        \While{nemixuje sa cesto pre upečenie všetkých chlebov}
        {
            Zahájenie mixer process\;
        }
        čakanie, až budú všetky chleby v bedničkách\;
        ukončenie Order process\;

        \caption{Zahájenie procesu výroby chleba}
        \label{algorithm:orderprocess}
    \end{algorithm}

    \begin{algorithm}[ht]
        \SetKw{Break}{break}
        Čakaj pokiaľ nieje daný proces hotový\;
        Spusť nasledujúci proces\;
        \caption{Proces výroby chleba}
        \label{algorithm:otherprocess}
    \end{algorithm}

    \subsection{Vstupy modelu}
    Vstupom simulačného modelu sú atribúty triedy \texttt{Args} v súbore \texttt{src/Args.h}:
    \begin{itemize}
        \item breads -- počet chlebov, ktoré se majú vyrobiť
        \item mixers -- počet mixérov
        \item tables -- počet stolov na krájanie
        \item carts -- počet vozíkov
        \item ovens -- počet pecí
        \item packers -- počet pracovníkov na balenie
    \end{itemize}

    Dalšie atribúty triedy \texttt{Args} sú pre ovládanie simulácie a výstupného súboru:
    \begin{itemize}
        \item outfile -- názov výstupného súboru
        \item simulations -- počet simulácií, ktoré se majú spustiť
    \end{itemize}

    Tieto vstupy sú zadávané pri spustení programu, viz.\ref{subsec:spousteni-simulacniho-model}

    \subsection{Výstup simulácie}
    Výstupom simulácie je možné riadit cez konštanty v súbore \texttt{src/macros.h}:
    \begin{itemize}
        \item SIMULAČNÝ REŽIM (DEBUG=0 a TEST=0) -- v tomto režimu je možné spustiť simuláciu a sledovať výsledky, model se spustí 3x.
        \item DEBUG REŽIM (DEBUG=1) -- v tomto režimu je možné sledovať podrobnejšie informácie o prebiehajúcich simuláciách
        v jednotlivých krokoch (doba vybavenia jednotlivých procesov, vstupy a výstupy jednotlivých procesov, \ldots)
        \item EXPERIMENT REŽIM (TEST=1) -- v tomto režimu je spustené experimentovanie pre zistenie,
        ktorá konfigurácia modelu (vstupu), je najvhôdnejšia pre pečenie
        chleba z hladiska nejmenšieho času vybavenia tzv. optimalnej konfigurácie.
    \end{itemize}

    Výstupom simulácie sú dáta o čase výroby chleba, zataženie jednotlivých strojov a pracovníkov.
    Výstup môže byť uložený do súboru, ktorý je zadaný príkazom \texttt{--outfile}, keď nie je zadaný
    výstup je vytlačený na štandratný výstup.

    \subsection{Spúšťanie simulačného modelu}\label{subsec:spousteni-simulacniho-model}
    Simulačný model se musí pred spustením skompilovať pomocou \texttt{Makefile}.
    \texttt{make} a spustený model je možné pomocou príkazu \texttt{make run} s predvolenými hodnotami vstupu.
    Model je možné spúšťať aj s parametrami, ktoré sú vstupmi simulačného modelu, viz. \texttt{src/Args.h}:
    Spúšťanie: \texttt{./build/bread\_factory --breads 100 --mixers 2 --tables 2
    --fermentations 2 --ovens 3 --loads 3 --simulations 3 --outfile out.txt}.


    \section{Podstata simulačných experimentov a ich priebeh}
    Cieľom práce, ktorá má simulovať výrobu chleba je zistiť, či pekáreň s ktorou pracujeme má najlepšiu konfiguráciu
    strojov, miestností a pracovníkov pre výrobu chleba.
    \subsection{Postup experimentovania}
    Pre experimentovanie je vytvorená funkcia \texttt{experiments()}, ktorá po spustení prebehne všetky možné
    konfigurácie počtu prostriedkov v rozmedzí 1 až 9.

    \subsection{Priebeh experimentov}
    Pri modelovaní výroby sme rozdelili prostriedky na konštanty a premenné.\\ \\
    \textbf{Konštanty:}
    \begin{itemize}
        \item Výstup jedného procesu mixovania -- \texttt{140kg cesta}
        \item Maximálna kapacita bochníkov vo vozíku -- \texttt{70 kusov}
        \item Maximálna kapacita vozíkov v peci -- \texttt{1 vozík}
    \end{itemize}
    \textbf{Premenné:}
    \begin{itemize}
        \item Počet chlebov na výrobu
        \item Počet mixérov
        \item Počet stolov/pracovníkov na krájanie
        \item Maximálna kapacita vozíkov vo fermentačnej miestnosti
        \item Počet pecí
        \item Počet pracovníkov na balenie
    \end{itemize}
    \textbf{Výsledky experimentov:}\\

    \subsubsection{Validácia modelu}
    \textbf{Reálna konfigurácia pekárne (simulácia)}
    \begin{center}
        \begin{tabular}{ |p{2cm}|p{2cm}|p{2cm}|p{2cm}|p{2cm}|p{2cm}|p{2cm}| }
            \hline
            čas / min & chleby / ks & mixéry & pracovníci na krájanie & kapacita fermentačnej miestnosti & kapacita pecí & pracovníci na balenie \\
            \hline\hline
            147 & 140 & 2 & 2 & 8 & 2 & 2  \\ \hline
            209 & 280 & 2 & 2 & 8 & 2 & 2 \\ \hline
            288 & 420 & 2 & 2 & 8 & 2 & 2 \\ \hline
        \end{tabular}
    \end{center}
    \textbf{Reálny konfigurácia pekárne (reálny získaný čas)}
    \begin{center}
        \begin{tabular}{ |p{2cm}|p{2cm}|p{2cm}|p{2cm}|p{2cm}|p{2cm}|p{2cm}| }
            \hline
            čas / min & chleby / ks & mixéry & pracovníci na krájanie & kapacita fermentačnej miestnosti & kapacita pecí & pracovníci na balenie \\
            \hline\hline
            150 & 140 & 2 & 2 & 8 & 2 & 2  \\ \hline
        \end{tabular}
    \end{center}
    Ako môžme vidieť reálny čas so simulovaným časom výroby chleba je totožný. Vďaka tomuto experimentu
    môžme pokračovať v experimentovaní s inými konfiguráciami, ktoré nám ukážu efektivitu pekárne.

    \subsubsection{Najlepšia konfigurácia experimenov (simulácia)}
    \begin{center}
        \begin{tabular}{ |p{2cm}|p{2cm}|p{2cm}|p{2cm}|p{2cm}|p{2cm}|p{2cm}| }
            \hline
            čas / min & chleby / ks & mixéry & pracovníci na krájanie & kapacita fermentačnej miestnosti & kapacita pecí & pracovníci na balenie \\
            \hline\hline
            112 & 140 & 2 & 4 & 6 & 8 & 6  \\ \hline
            125 & 280 & 3 & 4 & 9 & 5 & 9  \\ \hline
            131 & 420 & 4 & 6 & 6 & 8 & 9 \\ \hline
        \end{tabular}
    \end{center}


    \subsubsection{Najhoršia konfigurácia experimenov (simulácia):}
    \begin{center}
        \begin{tabular}{ |p{2cm}|p{2cm}|p{2cm}|p{2cm}|p{2cm}|p{2cm}|p{2cm}| }
            \hline
            čas / min & chleby / ks & mixéry & pracovníci na krájanie & kapacita fermentačnej miestnosti & kapacita pecí & pracovníci na balenie \\
            \hline\hline
            228 & 140 & 3 & 1 & 3 & 6 & 9  \\ \hline
            327 & 280 & 6 & 1 & 2 & 7 & 2 \\ \hline
            419 & 420 & 8 & 1 & 4 & 1 & 2 \\ \hline
        \end{tabular}
    \end{center}

    \subsection{Závery experimentov}
    Pri experimentovaní bolo spustených 177 147 konfigurácií, z ktorých sme zistili najlepšie a najhoršie nakonfigurovanie
    pekárne pre výrobu 140, 280 a 420 chlebov. Týmito vysledkami vieme tvrdit že pekáreň nemá najefektívnejšie nakonfigurovanie.
    Najlepší výsledok pre 140 chlebov bol o \textbf{38 minút} rýchlejší ako pri reálnom čase. Z priložených tabuliek
    môžme vidieť aj konfiguracie pre prípad výroby 280 a 420 chlebov. Do výsledkov sme taktiež zahrnuli najhoršie nakonfigurovanie
    pre prípad vyvarovania sa tychto konfigurácií.

    \section{Zhrnutie simulačných experimentov}
    Ako bolo uvedené v predchádzajúcej kapitole, bolo spustených 177 147 konfigurácií. Z týchto konfigurácií sme zistili
    najlepšie a najhoršie nakonfigurovanie pekárne pre výrobu 140, 280 a 420 chlebov.
    Týmito vysledkami vieme tvrdit že pekáreň nemá najefektívnejšie nakonfigurovanie.
    \newpage
    \bibliographystyle{czechiso}
    \bibliography{documentation}
\end{document}
